\documentclass[12pt,a4paper]{article}
\usepackage[utf8]{vietnam}
\usepackage{amsmath,amssymb}
\usepackage{tabularx}
\usepackage{tikz}
\usepackage{xcolor}
% Giảm lề tối đa tận dụng diện tích giấy (chuẩn máy in laser/photo)
\usepackage[top=0.6cm, bottom=0.6cm, left=0.6cm, right=0.6cm]{geometry}
\usepackage[dethi]{ex_test}

\setlength{\parskip}{0pt}

% Dòng kẻ chấm mờ đều đặn chuẩn viết tay (khoảng cách 7.2mm rộng rãi)
\newcommand{\dongke}[1]{%
  \begingroup%
  \setlength{\parskip}{0pt}%
  \baselineskip=20.5pt%
  \vspace{1.5mm}%
  \foreach \i in {1,...,#1}{%
    \noindent{\color{gray!60}\dotfill}\par%
  }%
  \vspace{1mm}%
  \endgroup%
}

% Định nghĩa ô bài tập tự luận trong bảng
\newcommand{\oBaiTap}[3][7]{%
  \vspace{1.5mm}%
  \noindent\textbf{#2} #3\par%
  \dongke{#1}%
}

% Câu trắc nghiệm 1 cột: 4 đáp án dàn hàng ngang trên 1 dòng
\newcommand{\cauTN}[6]{% \cauTN{số câu}{đề}{A}{B}{C}{D}
  \noindent\textbf{Câu #1.} #2\par\vspace{0.4mm}%
  \noindent%
  \makebox[0.25\linewidth][l]{\textbf{A.} #3}%
  \makebox[0.25\linewidth][l]{\textbf{B.} #4}%
  \makebox[0.25\linewidth][l]{\textbf{C.} #5}%
  \makebox[0.25\linewidth][l]{\textbf{D.} #6}\par\vspace{0.8mm}%
}

\begin{document}
\pagestyle{empty}

% #############################################################################
% ========================= TOÁN 6 - TUẦN 1 - BUỔI 2: MÃ ĐỀ 101 =============
% #############################################################################

% -------------------- TRANG 1 (MÃ ĐỀ 101): TRẮC NGHIỆM --------------------
\noindent
\begin{tabularx}{\linewidth}{|X|c|}
\hline
\textbf{BÀI 2: CÁCH GHI SỐ TỰ NHIÊN (TUẦN 1 - BUỔI 2)} & \textbf{Mã đề: 101} \\
\hline
\end{tabularx}

\vspace{1.2mm}

\noindent
\begin{tabularx}{\linewidth}{|X|}
\hline
\rule[-3mm]{0pt}{9.5mm}\textbf{Họ và tên:} \dotfill \quad \textbf{Lớp:} \makebox[1.8cm]{\dotfill} \\
\hline
\end{tabularx}

\vspace{1.2mm}

\noindent
\begin{tabularx}{\linewidth}{|p{3cm}|X|}
\hline
\centering\textbf{ĐIỂM SỐ} & \centering\textbf{LỜI NHẬN XÉT} \tabularnewline
\hline
\rule{0pt}{2.2cm} & \tabularnewline
\hline
\end{tabularx}

\vspace{1.5mm}
\noindent\textbf{BẢNG ĐIỀN ĐÁP ÁN TRẮC NGHIỆM (Mỗi câu đúng 0,25 điểm):}
\vspace{1.2mm}

\noindent
\renewcommand{\arraystretch}{1.5}
\begin{tabularx}{\linewidth}{|c|*{12}{>{\centering\arraybackslash}X|}}
\hline
\textbf{Câu} & \textbf{1} & \textbf{2} & \textbf{3} & \textbf{4} & \textbf{5} & \textbf{6} & \textbf{7} & \textbf{8} & \textbf{9} & \textbf{10} & \textbf{11} & \textbf{12} \\
\hline
\textbf{Đ/A} & & & & & & & & & & & & \\
\hline
\end{tabularx}
\renewcommand{\arraystretch}{1.0}

\vspace{3mm}

\noindent\textbf{\large PHẦN I. TRẮC NGHIỆM KHÁCH QUAN (3,0 điểm)}
\vspace{1.5mm}

\cauTN{1}{Số La Mã $\text{XIV}$ có giá trị trong hệ thập phân là:}{$16$}{$14$}{$15$}{$19$}

\cauTN{2}{Số $29$ được biểu diễn bằng chữ số La Mã là:}{$\text{XXIX}$}{$\text{XXVII}$}{$\text{XXXI}$}{$\text{XIX}$}

\cauTN{3}{Trong số $54\,823$, chữ số 4 có giá trị là:}{$400$}{$40$}{$40\,000$}{$4\,000$}

\cauTN{4}{Biểu diễn số $\overline{abc}$ trong hệ thập phân là:}{$a + b + c$}{$10a + b + c$}{$100a + 10b + c$}{$1000a + 100b + c$}

\cauTN{5}{Số liền trước của số $1\,000$ là:}{$1\,001$}{$999$}{$990$}{$900$}

\cauTN{6}{Số tự nhiên nhỏ nhất có ba chữ số khác nhau là:}{$100$}{$102$}{$123$}{$101$}

\cauTN{7}{Dùng ba chữ số $0; 3; 5$ viết được bao nhiêu số tự nhiên có hai chữ số khác nhau?}{$4$}{$6$}{$3$}{$5$}

\cauTN{8}{Số La Mã $\text{XXVI}$ có giá trị là:}{$24$}{$26$}{$27$}{$29$}

\cauTN{9}{Số liền sau của số tự nhiên lớn nhất có 3 chữ số là:}{$1\,000$}{$998$}{$1\,001$}{$999$}

\cauTN{10}{Biểu diễn số $\overline{ab}$ trong hệ thập phân là:}{$a + b$}{$10a + b$}{$a \cdot b$}{$10b + a$}

\cauTN{11}{Số tự nhiên lớn nhất có bốn chữ số khác nhau là:}{$9\,999$}{$9\,876$}{$9\,867$}{$9\,786$}

\cauTN{12}{Đồng hồ số La Mã có kim ngắn chỉ vào số $\text{IX}$ nghĩa là lúc:}{$11\text{ giờ}$}{$9\text{ giờ}$}{$4\text{ giờ}$}{$6\text{ giờ}$}

\newpage
% -------------------- TRANG 2 (MÃ ĐỀ 101): TỰ LUẬN --------------------
\noindent
\begin{tabularx}{\linewidth}{|X|c|}
\hline
\textbf{PHẦN II. TỰ LUẬN (7,0 ĐIỂM)} & \textbf{Mã đề: 101} \\
\hline
\end{tabularx}

\vspace{1.2mm}

\noindent
\begin{tabularx}{\linewidth}{|X|}
\hline
\rule[-3mm]{0pt}{9.5mm}\textbf{Họ và tên:} \dotfill \quad \textbf{Lớp:} \makebox[1.8cm]{\dotfill} \\
\hline
\end{tabularx}

\vspace{2.5mm}

\noindent\textbf{\large Câu 1 (2,5 điểm). Chuyển đổi hệ ghi số:}
\vspace{1.2mm}

\noindent
\begin{tabularx}{\linewidth}{|X|X|}
\hline
\oBaiTap[7]{a)}{Viết các số sau bằng chữ số La Mã: $16; 24; 28$. Viết giá trị hệ thập phân của các số La Mã: $\text{XVIII}; \text{XXV}; \text{XXIX}$.} & 
\oBaiTap[7]{b)}{Từ các chữ số $0; 2; 5; 7$, hãy viết số tự nhiên lớn nhất và số tự nhiên nhỏ nhất có bốn chữ số khác nhau.} \\
\hline
\end{tabularx}

\vspace{3.5mm}

\noindent\textbf{\large Câu 2 (2,5 điểm). Cấu tạo số và Tìm số:}
\vspace{1.2mm}

\noindent
\begin{tabularx}{\linewidth}{|X|X|}
\hline
\oBaiTap[7]{a)}{Phân tích các số sau thành tổng giá trị các chữ số của nó: $3\,746$ và $\overline{a8b}$.} & 
\oBaiTap[7]{b)}{Tìm một số tự nhiên có hai chữ số, biết rằng nếu viết thêm chữ số $2$ vào bên trái số đó thì được số mới gấp $6$ lần số ban đầu.} \\
\hline
\end{tabularx}

\vspace{3.5mm}

\noindent\textbf{\large Câu 3 (2,0 điểm). Bài toán thực tế:}
\vspace{1.2mm}

\noindent
\begin{tabularx}{\linewidth}{|X|}
\hline
\vspace{0.8mm}
Để đánh số trang một cuốn sách dày $120$ trang (từ trang $1$ đến trang $120$), người ta phải dùng tất cả bao nhiêu chữ số?
\par\vspace{1.5mm}
\noindent\textit{Bài giải:}
\dongke{8}
\\ \hline
\end{tabularx}

\newpage
% #############################################################################
% ========================= TOÁN 6 - TUẦN 1 - BUỔI 2: MÃ ĐỀ 102 =============
% #############################################################################

% -------------------- TRANG 3 (MÃ ĐỀ 102): TRẮC NGHIỆM --------------------
\noindent
\begin{tabularx}{\linewidth}{|X|c|}
\hline
\textbf{BÀI 2: CÁCH GHI SỐ TỰ NHIÊN (TUẦN 1 - BUỔI 2)} & \textbf{Mã đề: 102} \\
\hline
\end{tabularx}

\vspace{1.2mm}

\noindent
\begin{tabularx}{\linewidth}{|X|}
\hline
\rule[-3mm]{0pt}{9.5mm}\textbf{Họ và tên:} \dotfill \quad \textbf{Lớp:} \makebox[1.8cm]{\dotfill} \\
\hline
\end{tabularx}

\vspace{1.2mm}

\noindent
\begin{tabularx}{\linewidth}{|p{3cm}|X|}
\hline
\centering\textbf{ĐIỂM SỐ} & \centering\textbf{LỜI NHẬN XÉT} \tabularnewline
\hline
\rule{0pt}{2.2cm} & \tabularnewline
\hline
\end{tabularx}

\vspace{1.5mm}
\noindent\textbf{BẢNG ĐIỀN ĐÁP ÁN TRẮC NGHIỆM (Mỗi câu đúng 0,25 điểm):}
\vspace{1.2mm}

\noindent
\renewcommand{\arraystretch}{1.5}
\begin{tabularx}{\linewidth}{|c|*{12}{>{\centering\arraybackslash}X|}}
\hline
\textbf{Câu} & \textbf{1} & \textbf{2} & \textbf{3} & \textbf{4} & \textbf{5} & \textbf{6} & \textbf{7} & \textbf{8} & \textbf{9} & \textbf{10} & \textbf{11} & \textbf{12} \\
\hline
\textbf{Đ/A} & & & & & & & & & & & & \\
\hline
\end{tabularx}
\renewcommand{\arraystretch}{1.0}

\vspace{3mm}

\noindent\textbf{\large PHẦN I. TRẮC NGHIỆM KHÁCH QUAN (3,0 điểm)}
\vspace{1.5mm}

\cauTN{1}{Số La Mã $\text{XVI}$ có giá trị trong hệ thập phân là:}{$14$}{$16$}{$17$}{$19$}

\cauTN{2}{Số $24$ được biểu diễn bằng chữ số La Mã là:}{$\text{XXVI}$}{$\text{XXIV}$}{$\text{XXIX}$}{$\text{XIV}$}

\cauTN{3}{Trong số $73\,561$, chữ số 3 có giá trị là:}{$300$}{$3\,000$}{$30\,000$}{$30$}

\cauTN{4}{Biểu diễn số $\overline{xy}$ trong hệ thập phân là:}{$x + y$}{$x \cdot y$}{$10x + y$}{$10y + x$}

\cauTN{5}{Số liền sau của số $999$ là:}{$998$}{$1\,000$}{$1\,001$}{$1\,100$}

\cauTN{6}{Số tự nhiên lớn nhất có ba chữ số khác nhau là:}{$999$}{$987$}{$978$}{$897$}

\cauTN{7}{Dùng ba chữ số $0; 2; 7$ viết được bao nhiêu số tự nhiên có hai chữ số khác nhau?}{$4$}{$6$}{$3$}{$5$}

\cauTN{8}{Số La Mã $\text{XIX}$ có giá trị là:}{$21$}{$19$}{$18$}{$11$}

\cauTN{9}{Số liền trước của số tự nhiên nhỏ nhất có 3 chữ số là:}{$99$}{$100$}{$101$}{$90$}

\cauTN{10}{Biểu diễn số $\overline{mnp}$ trong hệ thập phân là:}{$100m + 10n + p$}{$m + n + p$}{$10m + n + p$}{$1000m + n + p$}

\cauTN{11}{Số tự nhiên nhỏ nhất có bốn chữ số khác nhau là:}{$1\,000$}{$1\,023$}{$1\,234$}{$1\,012$}

\cauTN{12}{Đồng hồ số La Mã có kim ngắn chỉ vào số $\text{IV}$ nghĩa là lúc:}{$6\text{ giờ}$}{$4\text{ giờ}$}{$9\text{ giờ}$}{$5\text{ giờ}$}

\newpage
% -------------------- TRANG 4 (MÃ ĐỀ 102): TỰ LUẬN --------------------
\noindent
\begin{tabularx}{\linewidth}{|X|c|}
\hline
\textbf{PHẦN II. TỰ LUẬN (7,0 ĐIỂM)} & \textbf{Mã đề: 102} \\
\hline
\end{tabularx}

\vspace{1.2mm}

\noindent
\begin{tabularx}{\linewidth}{|X|}
\hline
\rule[-3mm]{0pt}{9.5mm}\textbf{Họ và tên:} \dotfill \quad \textbf{Lớp:} \makebox[1.8cm]{\dotfill} \\
\hline
\end{tabularx}

\vspace{2.5mm}

\noindent\textbf{\large Câu 1 (2,5 điểm). Chuyển đổi hệ ghi số:}
\vspace{1.2mm}

\noindent
\begin{tabularx}{\linewidth}{|X|X|}
\hline
\oBaiTap[7]{a)}{Viết các số sau bằng chữ số La Mã: $17; 26; 29$. Viết giá trị hệ thập phân của các số La Mã: $\text{XIV}; \text{XXIV}; \text{XXVII}$.} & 
\oBaiTap[7]{b)}{Từ các chữ số $0; 1; 4; 8$, hãy viết số tự nhiên lớn nhất và số tự nhiên nhỏ nhất có bốn chữ số khác nhau.} \\
\hline
\end{tabularx}

\vspace{3.5mm}

\noindent\textbf{\large Câu 2 (2,5 điểm). Cấu tạo số và Tìm số:}
\vspace{1.2mm}

\noindent
\begin{tabularx}{\linewidth}{|X|X|}
\hline
\oBaiTap[7]{a)}{Phân tích các số sau thành tổng giá trị các chữ số của nó: $4\,582$ và $\overline{x5y}$.} & 
\oBaiTap[7]{b)}{Tìm một số tự nhiên có hai chữ số, biết rằng nếu viết thêm chữ số $3$ vào bên trái số đó thì được số mới gấp $5$ lần số ban đầu.} \\
\hline
\end{tabularx}

\vspace{3.5mm}

\noindent\textbf{\large Câu 3 (2,0 điểm). Bài toán thực tế:}
\vspace{1.2mm}

\noindent
\begin{tabularx}{\linewidth}{|X|}
\hline
\vspace{0.8mm}
Để đánh số trang một cuốn truyện dày $115$ trang (từ trang $1$ đến trang $115$), người ta phải dùng tất cả bao nhiêu chữ số?
\par\vspace{1.5mm}
\noindent\textit{Bài giải:}
\dongke{8}
\\ \hline
\end{tabularx}

\end{document}
